\subsection*{一步一步建立直观}

\paragraph{1. 让矩阵成为局部速度的生成规则}

对线性系统 $x'=Ax$，矩阵把当前位置变成速度，而不是直接把当前位置变成下一时刻的位置。右图用旋转生成矩阵演示这一点：小箭头是 $Ax$，整段时间后的状态则由 $e^{tA}x(0)$ 给出。二者一个是瞬时规则，一个是累计演化。



\paragraph{2. 从很小的时间步想象指数}

很短时间内，可以用 $(I+hA)x$ 近似新位置；不断缩小时间步并重复作用，便引出矩阵指数。加入负的径向项时，每步既旋转又收缩。本图使用这一特殊矩阵的精确闭式解，避免把数值步长造成的误差误认成方程本身的性质。



\paragraph{3. 特征方向解释增长，不能包办所有矩阵}

把径向项改正，模长随时间增长。当前矩阵由伸缩与平面旋转组成，结构特别简单；一般矩阵可能有剪切、不可对角化等情况，需要利用级数、分块或其他方法。本节 Lean 示例也只覆盖正文声明的具体矩阵恒等式，不证明全部谱理论。


